\documentclass[11pt]{article}
\usepackage[margin=1in]{geometry}
\usepackage{booktabs}
\usepackage{enumitem}
\usepackage{hyperref}
\usepackage{amsmath}
\usepackage{longtable}

\title{CSC483/583 Project Report Details: Weighted-Cole Retrieval, Cross-Encoder Reranking, and Qwen Reranker}
\author{}
\date{}

\begin{document}
\maketitle

\section*{Scope of This Draft}

This document collects the report material for the system variant we plan to present as the final implementation. It intentionally focuses only on the selected pipeline:

\begin{verbatim}
python -m src.run_test_100 \
  --weighted-cole \
  --include-category \
  --cross-encoder \
  --cross-encoder-mode chunks_100 \
  --cross-encoder-query-mode natural_questions_avg
\end{verbatim}

The selected final-stage model experiment outside the local run is the Qwen reranker path implemented in:

\begin{verbatim}
src/colab_score_qwen_verifier_inputs.py
\end{verbatim}

The primary reported metric is Precision@1, also called Top-1 accuracy. The reason is that Jeopardy requires a single final answer, and in this project the final answer is the title of one Wikipedia page.

\section{Core Implementation: Indexing and Retrieval}

\subsection{Task Definition}

The project asks us to answer 100 Jeopardy clues whose answers are Wikipedia page titles. The system receives a Jeopardy category and clue, retrieves candidate Wikipedia pages, and ranks page titles. A prediction is correct if the top-ranked page title matches the gold answer title, after normalization.

\subsection{Data Preparation}

The raw Wikipedia collection contains approximately 280,000 pages across 80 files. Each page begins with a title inside double square brackets. Our pipeline parses the raw files into SQLite article records, cleans the article bodies, separates redirect pages from normal article pages, and builds retrieval indexes over page-level documents. This satisfies the project requirement that each Wikipedia page be indexed as a separate document rather than indexing each source file as one document.

The relevant processed data and index artifacts are:

\begin{itemize}[leftmargin=*]
  \item \texttt{data/processed/wiki\_articles.sqlite3}: parsed article records.
  \item \texttt{data/processed/wiki\_articles\_step1\_clean.sqlite3}: cleaned article bodies.
  \item \texttt{data/processed/wiki\_redirects.sqlite3}: redirect mapping database.
  \item \texttt{index/whoosh\_cole\_index}: main Cole-style Whoosh index.
  \item \texttt{index/whoosh\_cole\_redirect\_index}: redirect title index.
  \item \texttt{index/whoosh\_cole\_first\_sentence\_index}: lead sentence index.
  \item \texttt{index/whoosh\_cole\_first\_two\_sentences\_index}: first two sentences index.
  \item \texttt{index/whoosh\_cole\_first\_paragraph\_index}: first paragraph index.
  \item \texttt{index/dpr\_faiss}: dense DPR FAISS indexes.
\end{itemize}

\subsection{Wikipedia-Specific Issues}

Wikipedia content contains markup, redirects, page titles with aliases, long article bodies, lead sections that summarize entities, and many pages that mention clue terms without being the answer. The implementation addresses these issues as follows:

\begin{itemize}[leftmargin=*]
  \item Redirect pages are indexed separately, then mapped back to canonical article targets.
  \item Article bodies are cleaned before indexing so retrieval is not dominated by raw wiki markup.
  \item Lead-only indexes are built because the first sentence, first two sentences, and first paragraph often contain concise entity descriptions.
  \item Dense DPR FAISS indexes are built for several passage variants so exact lexical overlap is not the only retrieval signal.
  \item Generated natural questions are used as additional dense queries to bridge from Jeopardy wording to natural question wording.
\end{itemize}

\subsection{Query Construction}

For the selected run we use \texttt{--include-category}. Therefore the sparse retrieval query is constructed from:

\begin{verbatim}
category + " " + clue
\end{verbatim}

The selected query mode is the default \texttt{full}, so the full category-plus-clue text is used rather than selecting only named entities or nouns.

For dense DPR retrieval, the default query embedding text uses the category and clue:

\begin{verbatim}
The category is: {category}. The clue is: {clue}
\end{verbatim}

In addition, the system uses five Qwen-generated natural-language questions for each clue. These are encoded as DPR query embeddings and used in two dense retrieval components: an average over five natural-question retrievals and one retrieval from a concatenation of the five generated questions.

\subsection{Why Convert Clues into Natural Questions}

Jeopardy clues are not usually written like ordinary search questions. They are often fragments, indirect descriptions, category-dependent prompts, or wordplay. However, the dense retriever, cross-encoder, and Qwen reranker are closer to question-document relevance models: DPR is designed around natural questions and Wikipedia passages, MS MARCO-style cross-encoders are trained to compare a query with a passage, and Qwen3-Reranker-4B expects an instruction, query, and document.

For that reason, the system rewrites each clue into five natural-language questions before dense retrieval and reranking. The goal is not to add outside knowledge or reveal the answer. The goal is to express the same information need in a form that matches the training distribution of question-versus-passage models. This is especially useful when the category supplies the answer type or hidden constraint. For example, a category such as \texttt{CAPITAL CITY CHURCHES} changes the intended target from a church page to a capital city page, and the natural-question prompt is designed to preserve that constraint.

\subsection{Prompt and Template Inventory}

The following prompts and fixed templates were used in the selected pipeline and related failed experiments.

\textbf{Natural-question generation with Qwen.} The system prompt was:

\begin{verbatim}
You rewrite Jeopardy-style clues into natural search and QA questions.
Return valid JSON only.
\end{verbatim}

The user prompt was:

\begin{verbatim}
Write exactly 5 natural-language questions from the category and clue.
Each question must preserve the same intended answer as the original
category and clue.
The category may provide the topic, answer type, or special instruction.
The clue provides the hint, fact, or indirect cue used to identify
the same answer.
Ask questions with the same answer as the original clue.
Use only the information in the category and clue.
Do not reveal the answer in the question.

Category: {category}
Clue: {clue}

Return exactly this JSON shape and nothing else:
{"questions":["...","...","...","...","..."]}
\end{verbatim}

\textbf{Sparse retrieval query.} The selected first-stage sparse query was:

\begin{verbatim}
{category} {clue}
\end{verbatim}

\textbf{DPR dense query template.} The default dense query text was:

\begin{verbatim}
The category is: {category}. The clue is: {clue}
\end{verbatim}

The two natural-question DPR variants used either the average of five question embeddings or one embedding from the concatenation of the five generated questions.

\textbf{Cross-encoder query/document template.} For the selected reranking mode, the query side is each generated natural question. The document side is:

\begin{verbatim}
Title: {title}
Article: {cleaned_article_chunk}
\end{verbatim}

The article is split into 100-token chunks. The best chunk score is taken for each generated question, and the five question scores are averaged.

\textbf{Qwen chat verifier prompt.} The failed chat-verifier path used:

\begin{verbatim}
Question: {natural_question}
Article: {evidence}

=============

Does this article contain the answer to the question?
Answer with exactly one word: Yes or No.
\end{verbatim}

When the tokenizer supported it, the script appended \texttt{/no\_think} and disabled thinking in the chat template.

\textbf{Qwen3-Reranker-4B prompt.} The reranker instruction was:

\begin{verbatim}
Given a Jeopardy clue rewritten as a natural question, retrieve articles that
contain enough information to identify the answer.
\end{verbatim}

The reranker input pair was:

\begin{verbatim}
<Instruct>: {instruction}
<Query>: {natural_question}
<Document>: {evidence}
\end{verbatim}

The Qwen reranker chat wrapper instructed the model to judge whether the document satisfies the query and to answer only \texttt{yes} or \texttt{no}. The score is computed from the next-token logits for those two labels, not from free-form generated text.

\textbf{Paraphrase experiment prompt.} The failed sparse-retrieval paraphrase experiment used:

\begin{verbatim}
paraphrase: {category_and_clue_text} </s>
\end{verbatim}

\subsection{Weighted-Cole Retrieval Components}

The selected first-stage retriever is \texttt{--weighted-cole}. It combines sparse Whoosh scores and dense DPR FAISS scores. Each component retrieves candidates, normalizes its scores, and contributes to a weighted sum. The current active weights for the selected system are:

\begin{table}[h]
\centering
\begin{tabular}{l r}
\toprule
Component & Weight \\
\midrule
Title BM25 & 0.0 \\
Redirect title BM25 & 0.0 \\
Full cleaned body BM25 & 15.0 \\
First sentence BM25 & 0.0 \\
First two sentences BM25 & 0.0 \\
First paragraph BM25 & 0.0 \\
DPR FAISS, default category+clue query & 1.0 \\
DPR FAISS, average over five natural questions & 1.0 \\
DPR FAISS, concatenated natural questions & 1.0 \\
Year match & 0.0 \\
Quote match & 0.0 \\
Category first sentence semantic match & 0.0 \\
Category first two sentences semantic match & 0.0 \\
\bottomrule
\end{tabular}
\end{table}

The effective first-stage score is:

\[
\begin{aligned}
\text{score}(d) ={}& 15 \cdot \text{body}(d) \\
&+ 1 \cdot \text{faiss}(d) \\
&+ 1 \cdot \text{natural\_questions\_avg\_faiss}(d) \\
&+ 1 \cdot \text{natural\_questions\_concat\_faiss}(d)
\end{aligned}
\]

Sparse scores are normalized by component maximum score. Dense FAISS scores are min-max normalized. Candidate documents are merged by canonical article identity, using \texttt{source\_file} and \texttt{article\_index}.

\subsection{Cross-Encoder Reranking}

After first-stage retrieval, the selected system reranks the top 100 candidates with:

\begin{verbatim}
cross-encoder/ms-marco-MiniLM-L-12-v2
\end{verbatim}

The selected cross-encoder mode is \texttt{chunks\_100}. Each candidate article is hydrated with its cleaned body text, formatted with its title, split into 100-token chunks, and scored against each query. The selected cross-encoder query mode is \texttt{natural\_questions\_avg}, so the five generated natural questions are scored against each article chunk. For a candidate article, the system takes the best chunk score for each natural question and averages the five question scores.

The cross-encoder reranked list is then fused with the initial weighted-Cole score. Both the initial score and cross-encoder score are min-max normalized among the top 100 candidates. The current best fusion weights are:

\[
\text{fused}(d) = 1.0 \cdot \text{initial\_normalized}(d)
+ 7.0 \cdot \text{cross\_encoder\_normalized}(d)
\]

\subsection{Qwen3-Reranker-4B Final-Step Experiment}

The Qwen reranker experiment is implemented in:

\begin{verbatim}
src/colab_score_qwen_verifier_inputs.py
\end{verbatim}

The input file exported for Colab is:

\begin{verbatim}
data/processed/qwen_verifier_inputs.jsonl
\end{verbatim}

This file contains 1000 rows: 100 clues times the top 10 fused candidate articles. Each row contains the clue, category, gold answer, five generated natural questions, candidate title, candidate rank, retrieval scores, and the first 360 tokens of cleaned article evidence.

The Colab script scores each candidate with two Qwen-based signals:

\begin{itemize}[leftmargin=*]
  \item \texttt{Qwen/Qwen3-4B}: a chat-style Yes/No verifier.
  \item \texttt{Qwen/Qwen3-Reranker-4B}: a reranker-style Yes/No relevance scorer.
\end{itemize}

For the Qwen3-Reranker-4B path, each of the five natural questions is paired with the candidate evidence. The model computes a relevance probability from the logits of \texttt{yes} and \texttt{no}. The final candidate score is the average of the five per-question reranker probabilities. The script outputs:

\begin{itemize}[leftmargin=*]
  \item \texttt{qwen\_chat\_yes\_probability}
  \item \texttt{per\_question\_chat\_yes\_probability}
  \item \texttt{qwen\_reranker\_probability}
  \item \texttt{per\_question\_reranker\_probability}
\end{itemize}

The completed Qwen score file is:

\begin{verbatim}
data/processed/qwen_verifier_scores.jsonl
\end{verbatim}

The Qwen reranker is a final reranking step over the exported top 10 fused candidates. It can improve the final order of these 10 candidates, but it cannot recover a correct answer that was not already present in the exported top-10 set.

\section{Coding with LLMs}

LLMs were used as coding assistants rather than as a full replacement for system design. The system architecture was guided by the project requirements: parse Wikipedia pages, build page-level indexes, retrieve candidates for Jeopardy clues, rerank the top candidates, and evaluate with an answer-title match metric.

\subsection{How LLMs Helped}

LLMs were used for:

\begin{itemize}[leftmargin=*]
  \item Explaining design alternatives for retrieval and reranking.
  \item Implementing scripts for exporting candidate rows for offline/Colab scoring.
  \item Debugging cache behavior and checking which pipeline stages were active.
  \item Drafting and refining report text and experiment descriptions.
  \item Comparing possible final-step models, including cross-encoder rerankers and Qwen rerankers.
\end{itemize}

\subsection{Prompting Style}

Prompts were usually concrete and code-specific. Instead of asking the LLM to build the whole project at once, we asked for specific tasks such as:

\begin{itemize}[leftmargin=*]
  \item Inspect the current retrieval code and summarize which weighted components are active.
  \item Add an exporter that writes the top fused candidates and evidence text to JSONL.
  \item Update the Colab scorer to output both Qwen chat verifier scores and Qwen reranker scores.
  \item Explain which metric is most appropriate for the project report.
\end{itemize}

\subsection{What Worked}

LLM assistance worked well for navigating the codebase, explaining the interaction between retrieval and reranking stages, writing small standalone scripts, and turning experimental results into report-ready explanations. It was especially useful for identifying when a component existed in code but had zero weight in the active system.

\subsection{What Did Not Work}

LLM suggestions needed verification against the actual code and cached outputs. Some retrieval notes became stale after weights changed, so the final report must be based on the current code and measured run output, not earlier notes. The Qwen verifier/reranker path also requires Colab execution, so code existence alone is not a measured performance result.

\section{Measuring Performance}

\subsection{Primary Metric}

We use Precision@1, also called Top-1 accuracy, as the primary metric. This is the most direct metric for the task because the system must return one final Wikipedia page title as the answer. A question is correct if the first-ranked title matches the gold answer title after normalization.

\[
P@1 = \frac{\# \text{questions with correct answer at rank 1}}{\# \text{questions}}
\]

The project handout lists P@1, NDCG, and MRR as possible metrics and asks us to choose and justify one. P@1 is more appropriate than NDCG for our main claim because there is only one correct page title for each question. MRR would also be reasonable, but it emphasizes the rank of the correct answer throughout a ranked list. Since Jeopardy requires a single answer, P@1 best matches the final system objective.

\subsection{Supporting Metrics}

We also report Precision@5, Precision@10, Precision@20, Precision@50, and Precision@100 as supporting metrics. These are not the primary metric, but they help analyze whether the correct answer is present in the candidate set before or after reranking. This is useful for error analysis because it distinguishes retrieval failures from final ranking failures.

\subsection{Measured Results}

The following result was produced by the selected cached run:

\begin{verbatim}
python -m src.run_test_100 \
  --weighted-cole \
  --include-category \
  --cross-encoder \
  --cross-encoder-mode chunks_100 \
  --cross-encoder-query-mode natural_questions_avg
\end{verbatim}

The run used 100 Jeopardy questions. Cache status confirmed that the final ranking and stage-2 cross-encoder results were loaded from cache:

\begin{verbatim}
Final ranking cache hits: 100
Stage-2 ranking cache hits: 100
\end{verbatim}

\begin{table}[h]
\centering
\small
\begin{tabular}{lccc}
\toprule
System stage & P@1 & P@5 & P@10 \\
\midrule
Plain Whoosh body-only & 0.0600 & 0.1400 & 0.1700 \\
Cole body-only, no category & 0.2100 & 0.4200 & 0.4900 \\
Cole body-only + category & 0.2200 & 0.4600 & 0.5300 \\
Weighted-Cole + category & 0.3100 & 0.5000 & 0.5900 \\
Cross-encoder only & 0.4000 & 0.6200 & 0.6500 \\
Fused initial + cross-encoder & 0.4100 & 0.6100 & 0.6700 \\
Qwen chat verifier & 0.0400 & 0.2900 & 0.5900 \\
Qwen3-Reranker-4B & \textbf{0.4700} & 0.5800 & 0.5900 \\
\bottomrule
\end{tabular}
\caption{Main stage comparison. P@1 is the primary metric; P@5 and P@10 are supporting metrics.}
\end{table}

The primary final system result is therefore:

\[
\boxed{P@1 = 0.4700}
\]

This means that the final Qwen3-Reranker-4B stage returns the correct Wikipedia title as the top answer for 47 out of 100 questions. The strongest fully local stage before Colab scoring is the fused initial + cross-encoder stage with P@1 = 0.4100.

\subsection{Full Local Reranking Metrics}

For the local weighted-Cole and cross-encoder run, the evaluator also reports deeper candidate-set metrics:

\begin{table}[h]
\centering
\small
\begin{tabular}{lcccccc}
\toprule
System stage & P@1 & P@5 & P@10 & P@20 & P@50 & P@100 \\
\midrule
Weighted-Cole + category & 0.3100 & 0.5000 & 0.5900 & 0.6400 & 0.6900 & 0.7300 \\
Cross-encoder only & 0.4000 & 0.6200 & 0.6500 & 0.6800 & 0.6900 & 0.7300 \\
Fused initial + cross-encoder & 0.4100 & 0.6100 & 0.6700 & 0.6800 & 0.7000 & 0.7300 \\
\bottomrule
\end{tabular}
\caption{Deeper local metrics. Qwen reranking is excluded here because it only sees the top 10 exported candidates.}
\end{table}

\subsection{Qwen Final-Stage Results}

The completed Colab run printed:

\begin{verbatim}
Qwen chat verifier ranking accuracy:
  Top-1: 0.0400
  Top-5: 0.2900
  Top-10: 0.5900
Qwen reranker ranking accuracy:
  Top-1: 0.4700
  Top-5: 0.5800
  Top-10: 0.5900
Evaluated questions: 100
\end{verbatim}

These values are reported exactly from \texttt{src/colab\_score\_qwen\_verifier\_inputs.py}. They use the scoring script's normalization and top-10 exported candidate set. We do not mix these numbers with a separate alias-aware local recomputation, because the goal is to report the output of the Qwen scoring experiment as it was run.

\section{Error Analysis}

\subsection{Rank-Bucket Movement}

The error analysis is based on where the correct page appears in the ranked list at each major stage. Rank 1 means the question is answered correctly under the primary P@1 metric. Ranks 2--100 indicate ranking errors that a later reranker may still fix. ``Below 100'' and ``Absent'' indicate retrieval failures for the top-100 cross-encoder stage.

\begin{table}[h]
\centering
\scriptsize
\begin{tabular}{lrrrrrrrr}
\toprule
Stage & Rank 1 & 2--5 & 6--10 & 11--20 & 21--50 & 51--100 & Below 100 & Absent \\
\midrule
Vanilla body-only Whoosh & 6 & 8 & 3 & 1 & 1 & 0 & 3 & 78 \\
Weighted-Cole + category & 31 & 19 & 9 & 5 & 5 & 4 & 18 & 9 \\
Cross-encoder only & 40 & 22 & 3 & 3 & 1 & 4 & 0 & 27 \\
Fused cross-encoder & 41 & 20 & 6 & 1 & 2 & 3 & 0 & 27 \\
Qwen3-Reranker-4B & 47 & \multicolumn{6}{c}{top-10 reranker; see Qwen metric table} & top-10 limited \\
\bottomrule
\end{tabular}
\caption{Measured rank buckets over the 100 questions. Qwen3-Reranker-4B is evaluated only on the exported top-10 fused candidates.}
\end{table}

The biggest early improvement is from vanilla retrieval to weighted-Cole. Vanilla body-only Whoosh places the answer at rank 1 only 6 times and fails to retrieve the gold page for 78 questions. Weighted-Cole + category raises rank-1 answers to 31 and places the correct page somewhere in the top 100 for 73 questions. This shows that most of the improvement comes from better retrieval and candidate generation, not only from final reranking.

The cross-encoder then attacks ranking errors within the top 100. It improves rank-1 answers from 31 to 40 while keeping top-100 recall fixed at 73. The fused score improves one more question at rank 1 and slightly improves P@10 and P@50. Finally, Qwen3-Reranker-4B reranks only the top-10 exported fused candidates and improves P@1 from 0.4100 to 0.4700.

\subsection{Transition Counts}

\begin{table}[h]
\centering
\scriptsize
\begin{tabular}{lrrrr}
\toprule
Transition & Fixed to Rank 1 & Lost Rank 1 & Stayed Rank 1 & Still not Rank 1 \\
\midrule
Vanilla body-only $\rightarrow$ Weighted-Cole & 27 & 2 & 4 & 67 \\
Weighted-Cole $\rightarrow$ Cross-encoder & 15 & 6 & 25 & 54 \\
Cross-encoder $\rightarrow$ Fused & 3 & 2 & 38 & 57 \\
Fused $\rightarrow$ Qwen3-Reranker-4B & 6 & -- & -- & -- \\
\bottomrule
\end{tabular}
\caption{Top-1 movement between stages. For the final Qwen transition, the official script result improves P@1 from 0.4100 to 0.4700, i.e., six additional rank-1 answers.}
\end{table}

The weighted hybrid retrieval stage is the largest early jump: it fixes 27 questions that vanilla body-only retrieval did not answer at rank 1. The cross-encoder fixes another 15, but also loses 6 previous rank-1 answers because it can over-promote semantically related but wrong pages. Fusion is a stabilizer: it preserves most cross-encoder wins and recovers a few cases using the original weighted-Cole score. Qwen3-Reranker-4B adds the final semantic judgment over top-10 candidates and gives the best P@1.

\subsection{Category-Level Error Movement}

The 100 questions are distributed across 30 Jeopardy categories, so many category groups are small. For that reason, the table below reports counts rather than percentages. Each cell after the vanilla column gives the number of questions answered correctly at rank 1 in that category, followed by the change from the previous stage. Qwen counts use the official Qwen scoring script's strict title normalization over the exported top-10 candidates.

\begin{table}[h]
\centering
\scriptsize
\begin{tabular}{lrrrrrr}
\toprule
Super-category & N & Vanilla & Weighted & Cross-enc. & Fused & Qwen rerank \\
\midrule
People / Biography & 19 & 2 & 9 (+7) & 10 (+1) & 11 (+1) & 14 (+3) \\
Geography / Places & 17 & 0 & 4 (+4) & 4 (+0) & 5 (+1) & 8 (+3) \\
History / Politics & 25 & 0 & 4 (+4) & 10 (+6) & 8 (-2) & 9 (+1) \\
Entertainment / Pop Culture & 14 & 1 & 6 (+5) & 7 (+1) & 7 (+0) & 5 (-2) \\
Business / Companies & 5 & 0 & 2 (+2) & 0 (-2) & 0 (+0) & 2 (+2) \\
Media / Journalism & 5 & 1 & 2 (+1) & 2 (+0) & 2 (+0) & 3 (+1) \\
Organizations / Institutions & 3 & 1 & 2 (+1) & 2 (+0) & 2 (+0) & 2 (+0) \\
Language / Wordplay & 3 & 0 & 0 (+0) & 0 (+0) & 0 (+0) & 1 (+1) \\
Mixed / General Knowledge & 9 & 1 & 2 (+1) & 5 (+3) & 6 (+1) & 3 (-3) \\
\midrule
\textbf{Total} & \textbf{100} & \textbf{6} & \textbf{31 (+25)} & \textbf{40 (+9)} & \textbf{41 (+1)} & \textbf{47 (+6)} \\
\bottomrule
\end{tabular}
\caption{Top-1 correct counts after grouping the 30 Jeopardy categories into 9 super-categories.}
\end{table}

This grouped view is easier to interpret than the full Table 5. Weighted-Cole helps most on \texttt{People / Biography}, \texttt{Entertainment / Pop Culture}, and \texttt{Geography / Places}. These are categories where the clue and category often point to a named entity, a role, or a location, and full-body BM25 plus dense retrieval are good at surfacing the right page into the candidate set.

The cross-encoder contributes its largest gain on \texttt{History / Politics}, where it adds six more rank-1 answers over weighted-Cole. That is consistent with clues that require aligning a question-like reformulation with evidence spread across a sentence or paragraph rather than relying on exact title or body-term overlap. The same stage is much less helpful for \texttt{Geography / Places} and can actively hurt \texttt{Business / Companies}, where many semantically related company and brand pages are strong distractors.

Qwen reranking helps most late on \texttt{People / Biography} and \texttt{Geography / Places}, adding three more rank-1 answers in each group. These are often cases where the correct answer is already in the top 10 and the evidence snippet contains a fairly direct identity statement. By contrast, \texttt{Entertainment / Pop Culture} and \texttt{Mixed / General Knowledge} get worse under Qwen reranking, which suggests that the final yes/no-style evidence judgment is less stable when clues are broad, comparative, lyrical, or depend on indirect association rather than a single explicit identification sentence.

Some super-categories remain hard for structural reasons. \texttt{Language / Wordplay} has very little lexical overlap with the answer page by design, so sparse retrieval and passage reranking have almost no direct evidence to exploit. \texttt{Business / Companies} is difficult because parent-company questions create near-miss distractors among subsidiaries, products, and industry pages. \texttt{History / Politics} improves strongly with the cross-encoder, but the fused stage gives back two wins, which shows that combining initial retrieval and semantic reranking is not always monotonic on categories with many closely related historical entities.

\subsection{Full Category Table}

The full per-category table is kept below for completeness.

\begin{center}
\scriptsize
\begin{longtable}{p{0.36\linewidth}rrrrrr}
\toprule
Category & N & Vanilla & Weighted & Cross-enc. & Fused & Qwen rerank \\
\midrule
\endfirsthead
\toprule
Category & N & Vanilla & Weighted & Cross-enc. & Fused & Qwen rerank \\
\midrule
\endhead
``TIN'' MEN & 4 & 1 & 1 (+0) & 2 (+1) & 3 (+1) & 2 (-1) \\
'80s NO.1 HITMAKERS & 5 & 1 & 1 (+0) & 2 (+1) & 2 (+0) & 0 (-2) \\
1920s NEWS FLASH! & 4 & 0 & 1 (+1) & 2 (+1) & 2 (+0) & 2 (+0) \\
AFRICAN CITIES & 4 & 0 & 1 (+1) & 1 (+0) & 1 (+0) & 3 (+2) \\
AFRICAN-AMERICAN WOMEN & 4 & 0 & 1 (+1) & 1 (+0) & 1 (+0) & 2 (+1) \\
BROADWAY LYRICS & 1 & 0 & 0 (+0) & 0 (+0) & 0 (+0) & 0 (+0) \\
CAMBODIAN HISTORY \& CULTURE & 3 & 0 & 0 (+0) & 1 (+1) & 1 (+0) & 1 (+0) \\
CAPITAL CITY CHURCHES & 3 & 0 & 0 (+0) & 0 (+0) & 0 (+0) & 1 (+1) \\
CEMETERIES & 2 & 0 & 1 (+1) & 1 (+0) & 1 (+0) & 1 (+0) \\
COMPLETE DOM-INATION & 3 & 0 & 0 (+0) & 0 (+0) & 0 (+0) & 1 (+1) \\
CONSERVATION & 3 & 0 & 0 (+0) & 1 (+1) & 1 (+0) & 1 (+0) \\
GOLDEN GLOBE WINNERS & 5 & 0 & 2 (+2) & 2 (+0) & 2 (+0) & 2 (+0) \\
GREEK FOOD \& DRINK & 3 & 0 & 2 (+2) & 2 (+0) & 3 (+1) & 3 (+0) \\
HE PLAYED A GUY NAMED JACK RYAN IN... & 3 & 0 & 3 (+3) & 3 (+0) & 3 (+0) & 3 (+0) \\
HISTORICAL HODGEPODGE & 4 & 0 & 1 (+1) & 2 (+1) & 1 (-1) & 2 (+1) \\
HISTORICAL QUOTES & 3 & 0 & 0 (+0) & 1 (+1) & 1 (+0) & 0 (-1) \\
I'M BURNIN' FOR YOU & 2 & 0 & 0 (+0) & 1 (+1) & 1 (+0) & 0 (-1) \\
NAME THE PARENT COMPANY & 5 & 0 & 2 (+2) & 0 (-2) & 0 (+0) & 2 (+2) \\
NEWSPAPERS & 5 & 1 & 2 (+1) & 2 (+0) & 2 (+0) & 3 (+1) \\
NOTES FROM THE CAMPAIGN TRAIL & 3 & 0 & 1 (+1) & 2 (+1) & 2 (+0) & 2 (+0) \\
OLD YEAR'S RESOLUTIONS & 4 & 0 & 1 (+1) & 1 (+0) & 0 (-1) & 1 (+1) \\
POETS \& POETRY & 3 & 0 & 0 (+0) & 0 (+0) & 1 (+1) & 1 (+0) \\
POTPOURRI & 3 & 0 & 1 (+1) & 2 (+1) & 2 (+0) & 1 (-1) \\
RANKS \& TITLES & 2 & 1 & 2 (+1) & 2 (+0) & 2 (+0) & 2 (+0) \\
SERVICE ORGANIZATIONS & 3 & 1 & 2 (+1) & 2 (+0) & 2 (+0) & 2 (+0) \\
STATE OF THE ART MUSEUM & 5 & 0 & 0 (+0) & 0 (+0) & 0 (+0) & 0 (+0) \\
THAT 20-AUGHTS SHOW & 1 & 0 & 0 (+0) & 0 (+0) & 0 (+0) & 0 (+0) \\
THE QUOTABLE KEATS & 2 & 0 & 0 (+0) & 0 (+0) & 0 (+0) & 2 (+2) \\
THE RESIDENTS & 3 & 0 & 2 (+2) & 2 (+0) & 2 (+0) & 3 (+1) \\
UCLA CELEBRITY ALUMNI & 5 & 1 & 4 (+3) & 5 (+1) & 5 (+0) & 4 (-1) \\
\midrule
\textbf{Total} & \textbf{100} & \textbf{6} & \textbf{31 (+25)} & \textbf{40 (+9)} & \textbf{41 (+1)} & \textbf{47 (+6)} \\
\bottomrule
\caption{Top-1 correct counts by category and stage. Category labels with long host instructions are shortened for readability.}
\end{longtable}
\end{center}

The category movement shows that most gains are not uniform. Weighted-Cole gives the largest improvements in categories with strong article-body lexical evidence or helpful answer-type context, such as \texttt{UCLA CELEBRITY ALUMNI}, \texttt{HE PLAYED A GUY NAMED JACK RYAN IN...}, \texttt{GREEK FOOD \& DRINK}, \texttt{GOLDEN GLOBE WINNERS}, \texttt{NAME THE PARENT COMPANY}, and \texttt{THE RESIDENTS}. In these categories, adding the category text and full-body Cole indexing helps the system retrieve the right entity pages before deeper reranking.

The cross-encoder helps categories where the correct page is already present but not ranked first, including \texttt{1920s NEWS FLASH!}, \texttt{CAMBODIAN HISTORY \& CULTURE}, \texttt{CONSERVATION}, \texttt{HISTORICAL QUOTES}, \texttt{I'M BURNIN' FOR YOU}, and \texttt{NOTES FROM THE CAMPAIGN TRAIL}. These categories often require matching a natural-language clue to a supporting sentence rather than just matching title terms. However, \texttt{NAME THE PARENT COMPANY} shows a failure mode: the cross-encoder lost two weighted-Cole wins, likely because company clues contain many semantically related brands, subsidiaries, and corporate pages that can look relevant without being the requested parent company.

The Qwen reranker is especially helpful for categories where final semantic judgment over top-10 evidence is enough to choose the right page, such as \texttt{AFRICAN CITIES}, \texttt{THE QUOTABLE KEATS}, \texttt{NAME THE PARENT COMPANY}, \texttt{THE RESIDENTS}, and \texttt{NEWSPAPERS}. It does not solve categories where the correct answer is missing from the exported top 10 or where the evidence snippet is not enough. \texttt{STATE OF THE ART MUSEUM} remains unsolved across all stages because those clues ask for a state from a museum name; the retrieved evidence often focuses on museum pages, locations, or broad art topics rather than directly ranking the state title. \texttt{BROADWAY LYRICS} and \texttt{THAT 20-AUGHTS SHOW} have only one example each, so their zero counts should be read as individual hard cases rather than category-level trends.

\subsection{Why Some Questions Are Answered Correctly}

Many correct questions can be handled because their clues contain strong lexical or semantic evidence for the answer page. Examples include questions with distinctive names, specific dates, quoted phrases, organization names, locations, or entity descriptions that appear in the Wikipedia article body. The category also helps in the weighted-Cole system when the clue is short or ambiguous, because it gives additional context such as \texttt{NEWSPAPERS}, \texttt{GOLDEN GLOBE WINNERS}, or \texttt{CEMETERIES}.

Dense DPR retrieval and generated natural questions help when the clue is phrased differently from the Wikipedia article. The cross-encoder improves ranking by jointly reading the natural question and candidate article chunks. Qwen3-Reranker-4B further improves the top-10 ordering by using a stronger reranker trained to judge whether a document satisfies a query.

\subsection{Error Classes}

The observed errors can be grouped into the following classes:

\begin{enumerate}[leftmargin=*]
  \item \textbf{Retrieval miss.} The correct page is not in the top 100 after weighted-Cole, so the cross-encoder cannot recover it. The fused system has 27 such top-100 failures.
  \item \textbf{Ranking miss.} The correct page is retrieved but not placed at rank 1. These are the cases targeted by the cross-encoder, fusion, and Qwen reranker.
  \item \textbf{Broad-topic distractor.} Broad pages match many clue terms but are not the answer. For example, a general media page can match newspaper clues while not being the specific newspaper answer.
  \item \textbf{Jeopardy clue requires indirect reasoning.} Some clues require connecting multiple facts or interpreting wordplay rather than matching the article text directly.
  \item \textbf{Category bias.} Category text helps the tuned weighted-Cole system, but in simple body-only Whoosh it hurts badly because raw category words can dominate the query.
  \item \textbf{Generated-question error.} Natural questions sometimes preserve the wrong focus or miss the key entity, so dense retrieval and reranking optimize for the wrong interpretation.
  \item \textbf{Evidence truncation.} The Qwen final step sees only the exported evidence snippet for each candidate. If important evidence is later in the article, the reranker cannot use it.
  \item \textbf{Verifier calibration failure.} The Qwen chat Yes/No probability is a poor ranking signal even with a 4B model. It is a chat verifier, not a model trained specifically for reranking.
\end{enumerate}

\subsection{Components and Ideas That Did Not Help Enough}

Several implemented ideas were useful for exploration but did not become part of the final weighted configuration.

\begin{itemize}[leftmargin=*]
  \item \textbf{Plain Whoosh body-only retrieval is weak.} It reaches only P@1 = 0.0600 because it lacks the Cole cleaning/scoring setup and all dense retrieval signals.
  \item \textbf{Naively adding category to plain Whoosh hurts.} Plain Whoosh with category falls to P@1 = 0.0200. Raw category terms can dominate or distort the simple body query.
  \item \textbf{Cole body-only is the better simple IR baseline.} Cole body-only reaches P@1 = 0.2100 without category and 0.2200 with category. This is the better ``vanilla IR body'' comparison for the report.
  \item \textbf{Lead-only indexes do not beat full body.} First sentence only reaches P@1 = 0.0900, first two sentences reaches 0.1200, and first paragraph reaches 0.1500. Full body Cole with category reaches 0.2200, so lead-only evidence is too narrow as a standalone signal.
  \item \textbf{Title-only and redirect-only are not useful standalone rankers.} Title-only and redirect-only both produced P@1 = 0.0000 as standalone signals. They remain possible auxiliary evidence, but their final weights are zero.
  \item \textbf{Paraphrasing did not help enough for sparse retrieval.} The paraphrase model mostly changed word order or minor function words while preserving the important nouns and proper nouns. That does not add much new lexical evidence for TF-IDF/BM25-style retrieval. The current local environment also cannot reproduce the paraphrase run without downloading the T5 Paraphrase PAWS model.
  \item \textbf{Qwen chat verification failed as a final ranker.} The Qwen3-4B chat verifier produced P@1 = 0.0400. Its Yes/No probability is poorly calibrated for ranking article candidates.
  \item \textbf{Other indexes exist but are zero-weighted in the final system.} Qwen-summary FAISS, category semantic indexes, quote/year matches, and lead-only indexes are present in code. Most are final-weight zero because they did not improve the selected weighted-Cole configuration enough.
\end{itemize}

\subsection{How the Improved System Addresses Errors}

The stage-2 cross-encoder directly targets ranking errors among retrieved candidates. It improves P@1 from 0.3100 for the initial weighted-Cole ranking to 0.4100 after fusion. This suggests that many errors in the initial retrieval stage are not complete retrieval failures; rather, the correct answer is present but needs better reranking.

Qwen3-Reranker-4B further addresses ranking errors among the top fused candidates. It improves P@1 from 0.4100 to 0.4700. Since this final step only sees the exported top 10 fused candidates, it helps reranking errors but not failures where the correct answer is absent from the top 10.

\section{Improved Implementation}

The improved implementation goes beyond a basic sparse IR system in three ways.

\subsection{Weighted Hybrid Retrieval}

The base Whoosh retrieval is extended with a weighted hybrid retriever. The selected weighted-Cole system combines full-body BM25 with DPR dense retrieval over several article passage variants. This improves recall compared with relying only on exact sparse term matching.

\subsection{Generated Natural Questions}

Each Jeopardy clue is rewritten into five natural-language questions using Qwen. These questions are used as dense retrieval queries and as cross-encoder reranking queries. This improves compatibility between Jeopardy-style clues and ordinary Wikipedia prose.

\subsection{Cross-Encoder Stage-2 Reranking}

The system reranks the top 100 retrieved candidates with a cross-encoder. This is a major improvement over a standard bag-of-words system because the cross-encoder jointly reads the question and article chunk. The fusion of initial and cross-encoder scores gives the best local P@1 before Colab reranking:

\[
0.4100
\]

\subsection{Qwen3-Reranker-4B Final-Step Reranking}

The final-step experiment exports the top 10 fused candidates and scores them with Qwen3-Reranker-4B in Colab. This follows the project suggestion that a large language model can be used to rerank the top candidates. The scorer computes a probability that the article contains enough information to answer the natural question, then averages over five natural questions per clue. This is the final best measured stage:

\[
P@1 = 0.4700
\]

The command for Colab is:

\begin{verbatim}
!pip install -q "transformers>=4.51.0" accelerate torch

!python colab_score_qwen_verifier_inputs.py \
  --input qwen_verifier_inputs.jsonl \
  --output qwen_verifier_scores.jsonl \
  --chat-model-name Qwen/Qwen3-4B \
  --reranker-model-name Qwen/Qwen3-Reranker-4B \
  --device cuda \
  --batch-size 8 \
  --reranker-batch-size 2
\end{verbatim}

The Qwen3-Reranker-4B stage improves over the fused local system by 0.0600 absolute P@1, or six additional questions answered correctly at rank 1.

\section*{Files to Cite in the Final Report}

\begin{itemize}[leftmargin=*]
  \item \path{src/run_test_100.py}: evaluation driver, query construction, cross-encoder reranking, score fusion, metric printing.
  \item \path{src/search.py}: weighted-Cole retrieval implementation and component weights.
  \item \path{src/export_qwen_verifier_inputs.py}: export of top fused candidate evidence for Colab scoring.
  \item Qwen scorer script: \path{src/colab_score_qwen_verifier_inputs.py}. It implements the Qwen chat verifier and Qwen3-Reranker-4B scorer.
  \item \path{data/processed/qwen_verifier_inputs.jsonl}: exported top 10 fused candidates for 100 questions.
\end{itemize}

\end{document}
